\paragraph{window-\texorpdfstring{$6$}{6} 自由商障碍、boundary 中心残余与有限加法账本的不可能性}
在 window-\(6\) 的纤维退化多值性、boundary 中心直和结构以及有限秩加法账本障碍已经确立之后，还可把“全局自由对称解释失败”与“局域中心残余仍然存在”压缩为如下两条直接后果。

\begin{theorem}[自由对称解释与有限加法账本的联合不可能性]\label{thm:xi-time-part60ab3-free-symmetry-finite-ledger-impossible}
不存在任何三元组 \((H,L,\Phi)\) 同时满足下列条件：
\begin{enumerate}
  \item \(H\) 为有限群，并且对每个 \(m\ge 1\)，都存在一个 \(H\) 在
  \[
  \Omega_m:=\{0,1\}^m
  \]
  上的自由作用，使其轨道划分恰等于 \(\Fold_m^{\mathrm{bin}}\) 的纤维划分；
  \item \(L\) 为有限生成交换群；
  \item \(\Phi:(\NN_{>0},\times)\to(\ZZ\oplus L,+)\) 为单射幺半群同态。
\end{enumerate}
\leanverified{paper\_xi\_time\_part60ab3\_free\_symmetry\_finite\_ledger\_impossible}
\end{theorem}
\begin{proof}
若第 \(1\) 条成立，则取 \(m=6\) 即得到 \(\Fold_6^{\mathrm{bin}}\) 的纤维划分可由某个有限群的自由轨道划分实现。然而推论 \ref{cor:foldbin6-degeneracy-multivalued-nonfree} 已证明这在 window-\(6\) 上不可能，因为退化度直方图
\[
2:8,\qquad 3:4,\qquad 4:9
\]
具有多值轨道大小，不能来自任何自由有限群作用。

另一方面，第 \(2\) 条与第 \(3\) 条的组合又被定理 \ref{thm:killo-godel-compression-not-finite-rank-homologizable} 直接排除：不存在有限生成交换群 \(L\) 使
\[
\Phi:(\NN_{>0},\times)\to(\ZZ\oplus L,+)
\]
成为单射幺半群同态。故上述三元组不可能存在。证毕。
\end{proof}

\begin{proposition}[window-\texorpdfstring{$6$}{6} 的 boundary 中心残余]\label{prop:xi-time-part60ab3-window6-boundary-center-residual}
\leanverified{paper\_xi\_time\_part60ab3\_window6\_boundary\_center\_residual}
记
\[
X_6^{\mathrm{bdry}}=\{100001,100101,101001\}.
\]
则这三条 boundary 词都对应二点纤维 \(d_6^{\mathrm{bin}}(w)=2\)。对每个 \(w\in X_6^{\mathrm{bdry}}\)，记 \(\tau_w\) 为该二点纤维上的唯一非平凡对合。则
\[
C_{\mathrm{bdry}}
:=
\langle \tau_w:\ w\in X_6^{\mathrm{bdry}}\rangle
\cong
(\ZZ/2\ZZ)^3
\hookrightarrow
Z(G_6^{\mathrm{bin}}).
\]
并且在交换化映射下，
\[
C_{\mathrm{bdry}}\longrightarrow \bigl(G_6^{\mathrm{bin}}\bigr)^{\mathrm{ab}}
\]
的像构成一个规范三维坐标子格，亦即 boundary parity 的中心残余在 abelianization 中仍保持独立可辨。
\end{proposition}
\begin{proof}
推论 \ref{cor:fold-bin-boundary-center-z2} 已把三条 boundary 二点纤维上的 sheet-flip 识别为规范中心嵌入
\[
(\ZZ/2\ZZ)^3\hookrightarrow Z(G_6^{\mathrm{bin}}),
\]
其中三条生成元正对应于 \(X_6^{\mathrm{bdry}}\) 中的三个 \(w\)。这就给出第一条结论。

再由定理 \ref{thm:xi-window6-boundary-parity-canonical-direct-summand}，
\[
\bigl(G_6^{\mathrm{bin}}\bigr)^{\mathrm{ab}}
\cong
(\ZZ/2\ZZ)^3\oplus (\ZZ/2\ZZ)^{18},
\]
其中第一因子恰由 boundary 三条中心 sheet-flip 的像生成。故 \(C_{\mathrm{bdry}}\) 在交换化中的像彼此独立，并形成一个规范三维坐标子格。证毕。
\end{proof}

\noindent
由此，window-\(6\) 的多对一折叠既不能被解释为自由有限对称的纯商，也不能被有限秩加法账本完全线性化；但这并不意味着群论痕迹全部消失。恰相反，boundary 三条二点纤维仍在中心层留下一个可规范识别的 \((\ZZ/2\ZZ)^3\) 残余，并在 parity-charge 的交换化通道中保持三维独立坐标。

\endinput
